% Einheit 5 von 5 – Zusammengesetzte Figuren (bank/flaechen/e5.jsonl, zusammenbau v0.1)
%% TODO Zweigzeile Teil 1: was hier gelernt wird (Satz in Schülersprache aus dem Einheitstitel) – Einheit 5
\einheitenkopf[e5]{Einheit 5 von 5 · Zusammengesetzte Figuren}
\zweigzeile{ab Kl. 5, je nach Buch bis Kl. 6 · P10 oft}
\setcounter{aufgabe}{25}

%% TODO Titel: Ich-kann-Satz fehlt in der Bank (Platzhalter: Kettenname) – Nr. 26
%% TODO Anweisung: ein Satz über der ersten Teilaufgabe fehlt in der Bank – Nr. 26
\begin{aufgabe}{Zusammengesetzte Figuren}
%% TODO \rechenplatz{n}: Schrittzahl des Grundfalls fehlt in der Bank (nur bei fester Schreibform, 2.3 b)
\begin{teile}
\teil In welche Teilflächen zerlegst du die Figur? Benenne sie.

\viereck{(0,0)}{(6,0)}{(6,3)}{(2,3)}
\teil Ein L-förmiges Beet besteht aus zwei Rechtecken: $8$ m × $3$ m und $3$ m × $4$ m. Fläche? A = \leerfeld[m²]
\teil Eine Giebelwand: unten ein Rechteck, $8$ m breit und $3$ m hoch, darauf ein Dreieck mit derselben Grundseite und $2{,}5$ m Höhe. Fläche der Wand? A = \leerfeld[m²]
\teil Einem Rechteck $9$ m × $7$ m fehlt an einer Ecke ein Rechteck $4$ m × $3$ m. Fläche? A = \leerfeld[m²]
\teil Aus einer Holzplatte $70$ cm × $45$ cm wird ein Quadrat mit $15$ cm Seitenlänge ausgesägt. Restfläche? A = \leerfeld[cm²]
\teil Eine quadratische Tischplatte mit $90$ cm Seitenlänge hat in der Mitte ein rundes Loch mit $r = 5$ cm für den Sonnenschirm. Restfläche? Runde auf ganze cm². A = \leerfeld[cm²]
\teil Einem Rechteck $11$ m × $6$ m fehlt an einer Ecke ein Rechteck $5$ m × $2$ m. Umfang der Restfigur? u = \leerfeld[m]
\teil Eine Eisbahn besteht aus einem Rechteck $24$ m × $14$ m und an beiden kurzen Seiten je einem Halbkreis mit $14$ m Durchmesser. Berechne die Fläche der Eisbahn. Runde auf ganze m². \hfill (P10 2017 OS) \\ A = \leerfeld[m²]
\end{teile}
\end{aufgabe}

%% TODO Titel: Ich-kann-Satz fehlt in der Bank (Platzhalter: Kettenname) – Nr. 27
%% TODO Anweisung: ein Satz über der ersten Teilaufgabe fehlt in der Bank – Nr. 27
\begin{aufgabe}{Sachaufgabe mit Entscheidung (reicht die Platte?)}
\begin{teile}
\teil Für eine Tischplatte $1{,}4$ m × $0{,}75$ m hat Paul eine Holzplatte mit $1{,}2$ m² Fläche, die $0{,}8$ m breit ist. Reicht sie?
\end{teile}
\end{aufgabe}

%% TODO Titel: Ich-kann-Satz fehlt in der Bank (Platzhalter: Kettenname) – Nr. 28
%% TODO Anweisung: ein Satz über der ersten Teilaufgabe fehlt in der Bank – Nr. 28
\begin{aufgabe}{Verschnitt in Prozent (Vorrat, Niveau III)}
\begin{teile}
\teil Aus einem Brett $2$ m × $0{,}25$ m werden $6$ Stücke je $0{,}3$ m × $0{,}25$ m gesägt. Wie viel Prozent des Bretts sind Verschnitt?
\end{teile}
\end{aufgabe}

%% TODO Titel: Ich-kann-Satz fehlt in der Bank (Platzhalter: Kettenname) – Nr. 29
%% TODO Anweisung: ein Satz über der ersten Teilaufgabe fehlt in der Bank – Nr. 29
\begin{aufgabe}{Zusammengesetzte Figuren – Fehler finden, Begründen, Anwendung}
\begin{teile}
\teil Ein Rasen ist $15$ m × $8$ m groß, darin liegt ein runder Teich mit $r = 2$ m. Ina rechnet die Rasenfläche: \rechnung{A &= 15 \cdot 8 + \pi \cdot 2^2 \\ A &\approx 132{,}6 \text{ m}^2} Wo ist der Fehler?
\teil Eine L-Figur lässt sich auf zwei Arten in zwei Rechtecke zerlegen. Kommt beide Male dieselbe Fläche heraus? Begründe.
\teil Ein L-förmiges Zimmer besteht aus den Rechtecken $5$ m × $4$ m und $2{,}5$ m × $2$ m. Es bekommt Parkett für $32$ € je m²; für Verschnitt kauft man $5\,\%$ mehr. Was kostet das Parkett?
\end{teile}
\end{aufgabe}
